\documentclass[a4paper,12pt]{article}
\usepackage[utf8]{inputenc}
\usepackage[T2A]{fontenc}
\usepackage[russian]{babel}
\usepackage{amsmath}
\usepackage{array}
\usepackage{booktabs}
\usepackage{geometry}
\geometry{left=2cm, right=2cm, top=2cm, bottom=2cm}

\begin{document}

\begin{center}
\Large\textbf{Задача о коммивояжере --- генетический алгоритм}\\[4pt]
\large Вариант 8
\end{center}

\section*{Исходные данные}

Матрица расстояний между городами:

\begin{center}
\begin{tabular}{c|ccccc}
\toprule
 & 1 & 2 & 3 & 4 & 5 \\
\midrule
1 & 0  & 1  & 1  & 11 & 10 \\
2 & 1  & 0  & 5  & 1  & 9  \\
3 & 1  & 5  & 0  & 6  & 10 \\
4 & 11 & 1  & 6  & 0  & 7  \\
5 & 10 & 9  & 10 & 7  & 0  \\
\bottomrule
\end{tabular}
\end{center}

Параметры алгоритма: размер популяции $N = 4$, вероятность мутации $P_m = 0{,}01$, вероятность скрещивания $P_c = 0{,}9$.

Оператор скрещивания --- двухточечный (для перестановок).
Оператор мутации --- случайная перестановка двух генов.
Оператор отбора --- рулеточный с весом $w_i = f_{\max} + f_{\min} - f_i$.

%======================================================
\section*{Этап 1. Исходная популяция}
%======================================================

Случайным образом сформирована начальная популяция из 4~перестановок.

Вычисление целевой функции (сумма расстояний замкнутого маршрута):
\begin{align*}
f(42351) &= d_{42}+d_{23}+d_{35}+d_{51}+d_{14} = 1+5+10+10+11 = 37,\\
f(43152) &= d_{43}+d_{31}+d_{15}+d_{52}+d_{24} = 6+1+10+9+1 = 27,\\
f(42315) &= d_{42}+d_{23}+d_{31}+d_{15}+d_{54} = 1+5+1+10+7 = 24,\\
f(23451) &= d_{23}+d_{34}+d_{45}+d_{51}+d_{12} = 5+6+7+10+1 = 29.
\end{align*}

Вероятности участия в размножении: $w_i = f_{\max}+f_{\min}-f_i = 37+24-f_i$:

\begin{center}
\textbf{Таблица 1. Исходная популяция}\\[4pt]
\begin{tabular}{c|c|c|c}
\toprule
\textnumero\ строки & Код & Значение ЦФ & Вероятность \\
\midrule
1 & 42351 & 37 & 24/127 \\
2 & 43152 & 27 & 34/127 \\
3 & 42315 & 24 & 37/127 \\
4 & 23451 & 29 & 32/127 \\
\bottomrule
\end{tabular}
\end{center}

%======================================================
\section*{Итерация 1}
%======================================================

\subsection*{Отбор пар (рулетка)}
Рулеточным отбором выбраны пары: $(2, 3)$ и $(3, 2)$.

\subsection*{Скрещивание}

\textbf{Пара 1:} родители 43152 и 42315. Точки разрыва: 1 и 3.

Разбиение: $(4|31|52)$ и $(4|23|15)$.

Обмен средними сегментами: $(*|23|*)$ и $(*|31|*)$.

Заполнение потомка~1 $(*|23|*)$: начинаем со второго элемента фрагмента родителя~1 $[31]$, т.\,е. с~1.
Обход родителя~1 по кругу от~1: $1, 5, 2, 4, 3$. Пропускаем $\{2,3\}$. Остаток: $[1, 5, 4]$.
$$\text{Потомок 1} = [1] + [23] + [54] = \mathbf{12354}.$$
$$f(12354) = d_{12}+d_{23}+d_{35}+d_{54}+d_{41} = 1+5+10+7+11 = 34.$$

Заполнение потомка~2 $(*|31|*)$: начинаем со второго элемента фрагмента родителя~2 $[23]$, т.\,е. с~3.
Обход родителя~2 от~3: $3, 1, 5, 4, 2$. Пропускаем $\{3,1\}$. Остаток: $[5, 4, 2]$.
$$\text{Потомок 2} = [5] + [31] + [42] = \mathbf{53142}.$$
$$f(53142) = d_{53}+d_{31}+d_{14}+d_{42}+d_{25} = 10+1+11+1+9 = 32.$$

\textbf{Пара 2:} та же пара $(3,2)$ --- аналогичные потомки: 53142~(32), 12354~(34).

Мутация: $P_m=0{,}01$ --- ни для одного потомка не сработала.

\begin{center}
\textbf{Таблица 2. Результат скрещивания 1-го поколения}\\[4pt]
\begin{tabular}{c|c|c|c}
\toprule
\textnumero\ строки & Родители & Потомки & Значение ЦФ \\
\midrule
2 & $4|31|52$ & 12354 & 34 \\
3 & $4|23|15$ & 53142 & 32 \\
3 & $4|23|15$ & 53142 & 32 \\
2 & $4|31|52$ & 12354 & 34 \\
\bottomrule
\end{tabular}
\end{center}

\subsection*{Редукция}

Объединяем родителей и потомков, сортируем по ЦФ, оставляем лучших~4:

\begin{center}
\textbf{Таблица 3. Популяция 1-го поколения}\\[4pt]
\begin{tabular}{c|c|c}
\toprule
\textnumero\ строки & Код & Значение ЦФ \\
\midrule
1 (3) & 42315 & 24 \\
2 (2) & 43152 & 27 \\
3 (4) & 23451 & 29 \\
4 (н) & 53142 & 32 \\
\bottomrule
\end{tabular}
\end{center}

%======================================================
\section*{Итерация 2}
%======================================================

\subsection*{Вероятности отбора}

$f_{\max}=32,\ f_{\min}=24,\ w_i = 32+24-f_i = 56-f_i$:

\begin{center}
\begin{tabular}{c|c|c|c}
\toprule
\textnumero & Код & ЦФ & Вероятность \\
\midrule
1 & 42315 & 24 & 32/112 \\
2 & 43152 & 27 & 29/112 \\
3 & 23451 & 29 & 27/112 \\
4 & 53142 & 32 & 24/112 \\
\bottomrule
\end{tabular}
\end{center}

\subsection*{Отбор и скрещивание}
Рулеткой выбраны пары: $(1, 2)$ и $(2, 1)$. Точки разрыва: 3 и 4.

\textbf{Пара 1:} родители 42315 и 43152. Разбиение: $(423|1|5)$ и $(431|5|2)$.

Потомок~1 $(*\!*\!*|5|*)$: фрагмент родителя~1 --- $[1]$, начинаем с~1.
Обход: $1, 5, 4, 2, 3$. Пропускаем $\{5\}$. Остаток: $[1, 4, 2, 3]$.
$$\text{Потомок 1} = [1,4,2]+[5]+[3] = \mathbf{14253}.\quad f = 11+1+9+10+1 = 32.$$

Потомок~2 $(*\!*\!*|1|*)$: фрагмент родителя~2 --- $[5]$, начинаем с~5.
Обход: $5, 2, 4, 3, 1$. Пропускаем $\{1\}$. Остаток: $[5, 2, 4, 3]$.
$$\text{Потомок 2} = [5,2,4]+[1]+[3] = \mathbf{52413}.\quad f = 9+1+11+1+10 = 32.$$

\textbf{Пара 2:} та же пара --- аналогичные потомки.

\begin{center}
\textbf{Таблица 4. Результат скрещивания 2-го поколения}\\[4pt]
\begin{tabular}{c|c|c|c}
\toprule
\textnumero\ строки & Родители & Потомки & ЦФ \\
\midrule
1 & $423|1|5$ & 14253 & 32 \\
2 & $431|5|2$ & 52413 & 32 \\
2 & $431|5|2$ & 52413 & 32 \\
1 & $423|1|5$ & 14253 & 32 \\
\bottomrule
\end{tabular}
\end{center}

\subsection*{Редукция}

Все потомки имеют ЦФ~$=32$, что хуже всех родителей. Популяция не меняется:

\begin{center}
\textbf{Таблица 5. Популяция 2-го поколения}\\[4pt]
\begin{tabular}{c|c|c}
\toprule
\textnumero\ строки & Код & Значение ЦФ \\
\midrule
1 (1) & 42315 & 24 \\
2 (2) & 43152 & 27 \\
3 (3) & 23451 & 29 \\
4 (4) & 53142 & 32 \\
\bottomrule
\end{tabular}
\end{center}

%======================================================
\section*{Итерация 3}
%======================================================

\subsection*{Вероятности отбора}

Те же, что в итерации~2: $w = [32, 29, 27, 24]$, $\sum w = 112$.

\subsection*{Отбор и скрещивание}

Рулеткой выбраны пары: $(3, 2)$ и $(1, 3)$.

\textbf{Пара 1:} родители 23451 и 43152. Точки разрыва: 2 и 3. Разбиение: $(23|4|51)$ и $(43|1|52)$.

Потомок~1 $(*\!*|1|*\!*)$: фрагмент родителя~1 --- $[4]$, начинаем с~4.
Обход родителя~1 от~4: $4, 5, 1, 2, 3$. Пропускаем $\{1\}$. Остаток: $[4, 5, 2, 3]$.
$$\text{Потомок 1} = [4,5]+[1]+[2,3] = \mathbf{45123}.\quad f = 7+10+1+5+6 = 29.$$

Потомок~2 $(*\!*|4|*\!*)$: фрагмент родителя~2 --- $[1]$, начинаем с~1.
Обход родителя~2 от~1: $1, 5, 2, 4, 3$. Пропускаем $\{4\}$. Остаток: $[1, 5, 2, 3]$.
$$\text{Потомок 2} = [1,5]+[4]+[2,3] = \mathbf{15423}.\quad f = 10+7+1+5+1 = 24.$$

\textbf{Пара 2:} родители 42315 и 23451. Случайное число $> P_c = 0{,}9$ --- скрещивание \textbf{не применяется}. Потомки --- копии родителей:
$$\text{Потомок 3} = 42315\ (24),\quad \text{Потомок 4} = 23451\ (29).$$

\begin{center}
\textbf{Таблица 6. Результат скрещивания 3-го поколения}\\[4pt]
\begin{tabular}{c|c|c|c}
\toprule
\textnumero\ строки & Родители & Потомки & ЦФ \\
\midrule
3 & $23|4|51$ & 45123 & 29 \\
2 & $43|1|52$ & 15423 & 24 \\
1 & 42315 (без скрещ.) & 42315 & 24 \\
3 & 23451 (без скрещ.) & 23451 & 29 \\
\bottomrule
\end{tabular}
\end{center}

\subsection*{Редукция}

Сортируем все 8 особей: $24, 24, 24, 27, 29, 29, 29, 32$. Лучшие~4:

\begin{center}
\textbf{Таблица 7. Популяция 3-го поколения}\\[4pt]
\begin{tabular}{c|c|c}
\toprule
\textnumero\ строки & Код & Значение ЦФ \\
\midrule
1 (1) & 42315 & 24 \\
2 (н) & 15423 & 24 \\
3 (н) & 42315 & 24 \\
4 (2) & 43152 & 27 \\
\bottomrule
\end{tabular}
\end{center}

%======================================================
\section*{Вывод}
%======================================================

\begin{center}
\begin{tabular}{l|c|c}
\toprule
Показатель & Исходная & После 3-й итерации \\
\midrule
Лучшее значение ЦФ & 24 & 24 \\
Среднее значение ЦФ & 29{,}25 & 24{,}75 \\
Худшее значение ЦФ & 37 & 27 \\
\bottomrule
\end{tabular}
\end{center}

За 3~итерации среднее значение целевой функции снизилось с $29{,}25$ до $24{,}75$.
Худшая особь улучшилась с $37$ до $27$.
Популяция сошлась к лучшему маршруту $4 \to 2 \to 3 \to 1 \to 5$ со стоимостью~$24$.

\end{document}
